% Post-CP17 candidate. This reduces the variation domain but does not prove sharpness.
\subsection{Two zero trace directions and the remaining matrix alternative}

\begin{lemma}[The two toral trace directions vanish]
\label{lem:DE-A3-positive-trace-radical}
\status{Candidate; complete proof written, independent audit pending.}
Keep the hypotheses of
Lemma~\ref{lem:DE-A3-positive-index-double-corank-candidate}.
Write $s=s_0+s_1$ for its two toral weight components and put
$R=\langle s_0,s_1\rangle=\langle s,es\rangle$.
The descended first-variation map annihilates $R$. Consequently
the quadratic determinant in that lemma is the determinant of
a linear map
\[
       H^0(B,H)/R\longrightarrow\operatorname{Hom}(k^2,k^2).
\]
This assertion does not prove that the determinant is nonzero.
\end{lemma}
\begin{proof}
We use the Cartier-dual description of the already proved
kernel--cokernel diagram. All vector spaces and maps in this
argument can be taken with their relative Frobenius twists.
Since the field is perfect, identifying those twists preserves
the kernels, cokernels, trace maps and zero identities used below.
Let $W=H^0(B,H)$, $\widetilde W=H^0(A,\pi^*H)$, and let
$C_s$ denote Cartier multiplication by $s$. The hypotheses give
\[
 \ker C_s=R,\qquad W=R\oplus\im C_s.
\]
Indeed Cartier is dual to the full coherent Frobenius, and its
induced map on $W/R$, dual to the Frobenius on $I_B$, is bijective.

The trace sequence is equivariant for the compatible toral
operators. Since the operator satisfies $e^2=e$, its two
eigenspace functors are exact. Choose $q_i\in\widetilde W_i$
with $\operatorname{Tr}q_i=s_i$, for $i=0,1$.
Cartier commutes with etale trace. The toral pairing identities
give $C_s(s_i)=0$, hence
\[
             C_{\pi^*s}(q_i)=\pi^*v_i
                 \quad\text{for some }v_i\in W.
\]
Choose $w_i\in W$ so that
$v_i+C_s(w_i)\in R$, and put $q_i'=q_i+\pi^*w_i$.
Thus $C_{\pi^*s}(q_i')\in\pi^*R$.
The classes of $q_0',q_1'$ in
$\widetilde W/\pi^*R$ form a basis of the kernel of its induced
Cartier map: they are independent because their traces are
$s_0,s_1$, and that kernel has dimension two by the preceding
lemma. They are therefore a basis of the dual of the cokernel
of the restricted coherent Frobenius.

Let $k\in K_B=\ker F_s$. By Serre duality, the matrix coefficient
of the first variation in a direction $t\in\widetilde W$,
tested against $\pi^*k$ and $q_j'$, is
\[
       \bigl\langle k,
          C_B\bigl(\operatorname{Tr}(tq_j')\,\omega_B\bigr)
       \bigr\rangle.
                                                        \tag{DE-TR1}
\]
The formula is simply the dual of restricting $F_t$ and then
passing to its cokernel; trace is dual to pullback. It also
follows by commuting relative Cartier with the etale trace.

Take $t=q_i$. The correction $\pi^*w_j$ in $q_j'$ contributes
$\beta(s_i,w_j)$ inside the pairing in (DE-TR1). This lies in
$\im C_s$: if $(w_j)_i$ is the weight-$i$ component, then
\[
                \beta(s_i,w_j)=\beta(s,(w_j)_i).
\]
It is annihilated by $k$. The remaining term is zero as well.
For $i\ne j$, the product $q_iq_j$ has toral weight one and
its Cartier image is zero. For $i=j$, the square identity
$C(q_i^2\omega_A)=q_iC(\omega_A)=0$ applies, since $A$ is mixed.
Thus (DE-TR1) vanishes for every $k$ and both $j$.
The descended variation in trace direction $s_i$ is zero.
The two directions span $R$, which proves the assertion.
\end{proof}

\begin{lemma}[The exact rank-one matrix alternative]
\label{lem:DE-A3-positive-matrix-ruling}
\status{Candidate; complete elementary proof written, independent audit pending.}
Let $k$ be an infinite field and $T:U\to\operatorname{Mat}_2(k)$
a linear map. Its determinant is identically zero if and only if
either every matrix in its image kills one fixed nonzero vector,
or all of their images lie in one fixed line.
In particular, for the preceding variation problem, nonvanishing
is equivalent to excluding both of these common annihilators.
Neither exclusion is established here.
\end{lemma}
\begin{proof}
The two stated alternatives imply the determinant vanishes.
Conversely, the zero matrix space satisfies either alternative.
Otherwise choose a nonzero image matrix. Since its determinant
is zero it has rank one, and changes of source and target bases
make it $A=\operatorname{diag}(1,0)$.
Write any other image matrix as $B=(b_{ij})$.
The identity $\det(A+tB)=0$ for all $t\in k$ gives $b_{22}=0$;
the identity $\det B=0$ then gives $b_{12}b_{21}=0$.
The image vector space is therefore the union of its two linear
subspaces defined by $b_{12}=0$ and $b_{21}=0$.
A vector space over an infinite field cannot be the union of
two proper linear subspaces. In the first case all matrices kill
the second basis vector; in the second case all images lie in
the first coordinate line. This proves the equivalence.
\end{proof}

\begin{remark}[Retained finite probe and its limitation]
\label{rem:DE-A3-positive-trace-probe-scope}
The program
\path{agents/obstructions_crystals/code/post_CP17_variation_weight_probe_v2.py}
checks one sharp characteristic-two product source. It gives
two toral weight spaces of dimension nine, variation-image
dimensions one and two on those spaces, and zero variation on
$s$ and $es$. Its predecessor and every earlier probe output are
retained. The v2 author run took $0.0163224049974815$ seconds.
These computations motivated the first lemma; its
proof is the trace and Cartier argument above. They do not
exclude either matrix alternative on an arbitrary mixed surface.
The earlier attempted factorization through the source Cartier
cokernel was refuted by the retained finite probe and is not
used: $W/R$ here is a different quotient.
\end{remark}
